% ── Appendix A — Detailed Methodology Tables ──
\chapter{Detailed Methodology Tables}
\label{app:detailed_methodology_tables}

This appendix contains the detailed methodology reference tables supporting
Chapter~\ref{ch:method}. It includes the raw dataset schema, Sentinel-2
Scene Classification Layer definitions, and the complete vegetation-index
formula reference used in the feature engineering stage.


% ============================================================
\section{Dataset Schema and Sentinel-2 Quality Classes}
% ============================================================

\small
\setlength{\LTpre}{6pt}
\setlength{\LTpost}{6pt}
\begin{longtable}{@{}p{2.6cm}p{1.4cm}p{1.6cm}p{7.4cm}@{}}
\caption{Dataset column descriptions, data types, and roles.}
\label{tab:columns}\\
\toprule
\textbf{Column} & \textbf{Type} & \textbf{Role} & \textbf{Description} \\
\midrule
\endfirsthead
\caption[]{Dataset column descriptions, data types, and roles (continued).}\\
\toprule
\textbf{Column} & \textbf{Type} & \textbf{Role} & \textbf{Description} \\
\midrule
\endhead
\midrule
\multicolumn{4}{r}{\textit{Continued on next page}}\\
\endfoot
\bottomrule
\endlastfoot
\texttt{LONGITUDE} & float64 & Metadata &
    WGS84 longitude of the pixel centroid. Used for spatial distribution
    plots; not a model input. \\
\addlinespace
\texttt{LATITUDE} & float64 & Metadata &
    WGS84 latitude of the pixel centroid. Used for zoom-level ground
    resolution calculation and spatial diagnostics. \\
\addlinespace
\texttt{B01--B12, B8A}
    & int32 & Spectral Signal &
    Sentinel-2 Level-2A spectral reflectance bands stored as scaled
    integers (Digital Numbers). Wavelengths span 443~nm (B01, coastal
    aerosol) to 2190~nm (B12, SWIR-2), covering the visible, red-edge,
    near-infrared, and shortwave infrared regions. These bands are the
    primary input to all vegetation index formulas and statistical feature
    extraction. Scaling to physical reflectance is described in
    Section~\ref{subsec:scaling}. \\
\addlinespace
\texttt{AOT} & int32 & Quality Flag &
    Aerosol Optical Thickness at 550~nm, encoded as a scaled integer.
    Used to identify and mask observations affected by elevated aerosol
    loading. \\
\addlinespace
\texttt{SCL} & int32 & Quality Flag &
    Sentinel-2 Scene Classification Layer. A per-pixel quality label
    generated by the Sen2Cor processor (see Table~\ref{tab:scl}).
    Used to identify and mask cloud, shadow, and other non-surface pixels. \\
\addlinespace
\texttt{CLP} & int32 & Quality Flag &
    Cloud Probability (0--255, 8-bit). Used to mask observations with
    elevated cloud contamination probability. \\
\addlinespace
\texttt{NDVI} & float64 & Derived &
    Normalised Difference Vegetation Index, pre-computed and included
    in the raw dataset. Values lie in $[-1, 1]$ by definition. \\
\addlinespace
\texttt{TRUECOLOR\_R/G/B} & int32 & Metadata &
    Red, green, and blue channels of an 8-bit true-colour composite.
    Used for visualisation only; not fed to any model. \\
\addlinespace
\texttt{QUADKEY} & str & Metadata &
    Bing Maps QUADKEY tile identifier at zoom level~22. The fundamental
    spatial unit of analysis throughout the pipeline; used for feature
    aggregation and train/test splitting. See Section~\ref{sec:quadkey}. \\
\addlinespace
\texttt{START\_DATE} & datetime64 & Metadata &
    Sentinel-2 acquisition date (start of 7-day compositing window).
    Parsed to \texttt{datetime64[ns]} and used for temporal filtering,
    weekly deduplication, and interpolation. \\
\addlinespace
\texttt{END\_DATE} & datetime64 & Metadata &
    End of the 7-day compositing window. Used for date validation. \\
\addlinespace
\texttt{landcover} & str & Target &
    Crop class label as collected in the field (raw string). Standardised
    to uppercase and whitespace-normalised during preprocessing. The
    classes retained for classification are determined in
    Section~\ref{subsec:label_std}. \\
\addlinespace
\texttt{wkt} & str & Metadata &
    Well-Known Text (WKT) geometry string of the source field polygon.
    Used for compactness quality control
    (Section~\ref{subsec:compactness}) and boundary pixel detection. \\
\addlinespace
\texttt{id} & int32 & Metadata &
    Numeric field identifier linking each QUADKEY to its source
    ground-truth polygon. Used as the grouping variable for field-level
    train/test splitting. \\
\addlinespace
\texttt{year} & int16 & Derived &
    ISO calendar year derived from \texttt{START\_DATE}. Used for
    weekly coverage diagnostics. \\
\addlinespace
\texttt{week} & int16 & Derived &
    ISO calendar week number derived from \texttt{START\_DATE}. Used
    for temporal deduplication. \\
\end{longtable}
\normalsize

\begin{table}[H]
\centering
\caption{Sentinel-2 Scene Classification Layer (SCL) class definitions.}
\label{tab:scl}
\begin{tabular}{@{}clp{8cm}@{}}
\toprule
\textbf{SCL Value} & \textbf{Class Name} & \textbf{Meaning} \\
\midrule
0  & No Data              & Missing or out-of-swath pixel \\
1  & Saturated / Defective & Sensor saturation or bad pixel \\
2  & Dark Area Pixels     & Topographic shadow or dark surface \\
3  & Cloud Shadows        & Shadow cast by clouds \\
4  & Vegetation           & Healthy vegetated surface \\
5  & Bare Soils           & Exposed soil or bare ground \\
6  & Water                & Open water surface \\
7  & Unclassified         & Could not be assigned to any class \\
8  & Cloud (medium prob.) & Medium-probability cloud \\
9  & Cloud (high prob.)   & High-probability cloud \\
10 & Thin Cirrus          & Semi-transparent high-altitude cloud \\
11 & Snow or Ice          & Snow or ice-covered surface \\
\bottomrule
\end{tabular}
\end{table}


% ============================================================
\section{Vegetation Index Formula Reference}
% ============================================================

\small
\setlength{\LTpre}{6pt}
\setlength{\LTpost}{6pt}
\begin{longtable}{@{}p{2.2cm}p{4.4cm}p{2.0cm}p{4.4cm}@{}}
\caption{Vegetation indices computed from scaled Sentinel-2 bands.
Band wavelength centres in nm are given for reference.
A small constant $\varepsilon = 10^{-10}$ is added to all denominators
to prevent division-by-zero errors.}
\label{tab:vi}\\
\toprule
\textbf{Index} & \textbf{Formula} & \textbf{Bands} & \textbf{Agronomic Signal} \\
\midrule
\endfirsthead
\caption[]{Vegetation indices computed from scaled Sentinel-2 bands (continued).}\\
\toprule
\textbf{Index} & \textbf{Formula} & \textbf{Bands} & \textbf{Agronomic Signal} \\
\midrule
\endhead
\midrule
\multicolumn{4}{r}{\textit{Continued on next page}}\\
\endfoot
\bottomrule
\endlastfoot
\multicolumn{4}{@{}l}{\textit{Greenness / Canopy Vigour}} \\
\addlinespace
NDVI &
$\dfrac{B08 - B04}{B08 + B04}$ &
B08 (842\,nm), B04 (665\,nm) &
General vegetation health and canopy density; the most widely used
crop monitoring index. \\
\addlinespace
SAVI &
$\dfrac{1.5(B08 - B04)}{B08 + B04 + 0.5}$ &
B08, B04 &
Soil-adjusted variant of NDVI; reduces bare-soil background effects
in sparsely vegetated fields~\cite{huete1988}. \\
\addlinespace
OSAVI &
$\dfrac{B08 - B04}{B08 + B04 + 0.16}$ &
B08, B04 &
Optimised soil-adjusted index; fixed soil adjustment factor
minimises sensitivity to soil brightness variation~\cite{rondeaux1996}. \\
\addlinespace
EVI &
$\dfrac{2.5(B08 - B04)}{B08 + 6 \cdot B04 - 7.5 \cdot B02 + 1}$, clipped to $[-1,1]$ &
B08, B04, B02 (490\,nm) &
Enhanced Vegetation Index; reduces atmospheric and canopy
background effects; sensitive to dense canopies where NDVI
saturates~\cite{huete2002}. \\
\addlinespace
GNDVI &
$\dfrac{B08 - B03}{B08 + B03}$ &
B08, B03 (560\,nm) &
Green NDVI; sensitive to chlorophyll concentration; complements
NDVI for crop nitrogen status monitoring~\cite{gitelson1996}. \\
\addlinespace
\multicolumn{4}{@{}l}{\textit{Water / Moisture Indices}} \\
\addlinespace
NDWI &
$\dfrac{B03 - B08}{B03 + B08}$ &
B03, B08 &
Normalised Difference Water Index; sensitive to open water and
canopy water content; useful for detecting waterlogged rice paddies. \\
\addlinespace
LSWI &
$\dfrac{B08 - B11}{B08 + B11}$ &
B08, B11 (1610\,nm) &
Land Surface Water Index; sensitive to leaf and soil moisture;
useful for irrigated crop detection~\cite{xiao2005}. \\
\addlinespace
NDMI &
$= \mathrm{LSWI}$ (same formula in Sentinel-2 context) &
B08, B11 &
Normalised Difference Moisture Index; equivalent to LSWI for
Sentinel-2 band configuration~\cite{wilson2002}. \\
\addlinespace
\multicolumn{4}{@{}l}{\textit{Burn / Dryness / Vegetation Stress Indices}} \\
\addlinespace
NBR &
$\dfrac{B08 - B12}{B08 + B12}$ &
B08, B12 (2190\,nm) &
Normalised Burn Ratio; sensitive to charred vegetation and dry
biomass; provides contrast for post-harvest field detection. \\
\addlinespace
NBR2 &
$\dfrac{B11 - B12}{B11 + B12}$ &
B11, B12 &
Second Normalised Burn Ratio; differentiates healthy vegetation from
stressed or dried biomass using two SWIR bands. \\
\addlinespace
\multicolumn{4}{@{}l}{\textit{Red-Edge / Chlorophyll-Sensitive Indices}} \\
\addlinespace
NDRE &
$\dfrac{B08 - B05}{B08 + B05}$ &
B08, B05 (705\,nm) &
Red-edge NDVI; highly sensitive to chlorophyll content and early
crop stress; advantageous for crop type discrimination~\cite{gitelson1994}. \\
\addlinespace
NDRE\_B06 &
$\dfrac{B08 - B06}{B08 + B06}$ &
B08, B06 (740\,nm) &
Red-edge variant using the second red-edge band; captures a
different stage of the red-edge reflectance slope. \\
\addlinespace
NDRE\_B07 &
$\dfrac{B08 - B07}{B08 + B07}$ &
B08, B07 (783\,nm) &
Red-edge variant using the third red-edge band; sensitive to
dense canopy chlorophyll. \\
\addlinespace
NDRE\_B8A\_B05 &
$\dfrac{B8A - B05}{B8A + B05}$ &
B8A (865\,nm), B05 &
Narrow-NIR red-edge index; increases sensitivity to canopy
structure at the red-edge inflection. \\
\addlinespace
CI\_RE &
$\dfrac{B08}{B05} - 1$ &
B08, B05 &
Chlorophyll Index Red-Edge; directly proportional to leaf
chlorophyll content; more linear than NDRE at high biomass
levels~\cite{gitelson2003}. \\
\end{longtable}
\normalsize
